\documentclass[11pt]{article}
\usepackage[T1]{fontenc}
\usepackage{lmodern}
\usepackage[margin=1in]{geometry}
\usepackage{amsmath,amssymb,amsthm,mathtools}
\usepackage{enumitem}
\usepackage[hidelinks]{hyperref}
\usepackage{microtype}
\hypersetup{
 pdftitle={Positive lower density of large prime gaps},
 pdfauthor={OpenAI},
 pdfsubject={Ordinary lower density of consecutive prime gaps}}
\pdfinfoomitdate=1
\pdftrailerid{}
\pdfsuppressptexinfo=15
\newtheorem{theorem}{Theorem}[section]
\newtheorem{proposition}[theorem]{Proposition}
\newtheorem{lemma}[theorem]{Lemma}
\newtheorem{corollary}[theorem]{Corollary}
\theoremstyle{remark}
\newtheorem{remark}[theorem]{Remark}
\DeclareMathOperator{\Var}{Var}
\title{Positive lower density of large prime gaps}
\author{OpenAI}
\date{September 25, 2026}
\begin{document}
\maketitle
\begin{abstract}
For every fixed $C>0$, we prove that a positive proportion of consecutive
prime gaps exceed $C\log p$, where $p$ is the smaller prime. The proportion
is bounded below for every sufficiently large initial segment of the
prime sequence, with a constant depending on $C$. It follows that the
indices at which $p_n/n$ increases have positive lower asymptotic density,
answering a question of Erd\H{o}s and Prachar.
\end{abstract}
\tableofcontents
\section{Introduction}

Let $p_n$ denote the $n$th prime and write $d_n=p_{n+1}-p_n$.
The prime number theorem identifies $\log p$ as the average scale of
prime gaps near $p$. Individual gaps can be much larger.
Westzynthius proved that $d_n/\log p_n$ is unbounded
\cite{Westzynthius1931}; the constructions of Erd\H{o}s and Rankin
gave increasingly strong quantitative lower bounds for exceptionally
large gaps \cite{Erdos1935,Rankin1938}.
Modern refinements were obtained independently by Ford, Green,
Konyagin and Tao and by Maynard
\cite{FordGreenKonyaginTao2016,Maynard2016Large}, followed by their
joint work \cite{FordGreenKonyaginMaynardTao2018}.
These results concern the size of very large gaps. A lower bound for
the largest gap does not by itself give a positive proportion of
gaps exceeding a fixed multiple of $\log p$.

We study this frequency question with equal weight on the prime
indices. For a set $A$ of positive integers, its lower asymptotic
density is $\liminf_{N\to\infty}|A\cap[1,N]|/N$.

\begin{theorem}\label{thm:main}
For every fixed real $C>0$, there are constants $c(C)>0$ and $N_0(C)$
such that, for every integer $N\ge N_0(C)$,
\[
 \#\{1\le n\le N:p_{n+1}-p_n>C\log p_n\}\ge c(C)N.
\]
\end{theorem}

The density is ordinary lower density: the bound holds for every
sufficiently large initial segment, and both constants may depend
on the fixed threshold $C$.

One consequence concerns successive values of $p_n/n$.
Erd\H{o}s and Prachar asked whether the indices at which this ratio
increases have positive lower density \cite[p.~256]{ErdosPrachar1962}.

\begin{corollary}\label{cor:ratios}
The set
\[
 \left\{n\ge1:\frac{p_n}{n}<\frac{p_{n+1}}{n+1}\right\}
\]
has positive lower asymptotic density.
\end{corollary}
\begin{proof}
The displayed inequality is equivalent to
$p_{n+1}-p_n>p_n/n$.
The prime number theorem gives $p_n/n\sim\log p_n$.
Apply Theorem~\ref{thm:main} with $C=2$ and discard the finitely many
indices for which $p_n/n\ge2\log p_n$.
\end{proof}

Thus the corollary answers the Erd\H{o}s--Prachar question affirmatively.

\subsection*{Gap frequency and empty intervals}

The distinction between counting intervals and counting prime gaps is
essential. A long gap contains many integer starting points of
prime-free short intervals. Consequently a positive proportion of
empty intervals may be accounted for by a sparse set of very long
gaps. Bazzanella, Languasco and Zaccagnini distinguish the two counting
measures explicitly \cite{BLZ2010}. Their Theorem~5 gives positive
prime-start proportions for each fixed threshold below
$2/3.454=0.579038\ldots$, whereas their Theorem~3 reaches some
intervals longer than the average prime spacing by counting integer
starts. Tao's empty-interval argument \cite{Tao2019} also identifies
the additional control required to pass from interval measure to
gap counts.

Conditional results describe much more of the distribution.
Gallagher derived Poisson statistics for prime counts in
integer-start intervals of length $\lambda\log X$, for fixed
$\lambda>0$, from a uniform
Hardy--Littlewood prime-tuple hypothesis
\cite[Theorem~1]{Gallagher1976}.
The Poisson-tail results of Jha
\cite[Theorems~1.4 and~1.6]{Jha2026} treat growing parameter ranges
under a stronger uniform tuple hypothesis.
The counting measures and uniformity assumptions matter in comparing
these conclusions with Theorem~\ref{thm:main}.

\subsection*{The adjacent-interval construction}

We control the number of distinct gaps by requiring a prime just
before an empty interval. Fix $\lambda>0$, put
$h=\lfloor\lambda\log X\rfloor$, and average over integers
$X<m\le2X$. We construct a nonnegative weight for which a prime in
$(m,m+h]$ has substantial weighted probability, whereas the expected
number of primes in $(m+h,m+2h]$ is arbitrarily small.
One second-moment estimate detects the first prime; a separate
second-moment estimate converts the resulting weighted mass into a
positive proportion of ordinary integer starts.
The last prime in $(m,m+h]$ then begins a gap longer than $h$.
It can be selected by at most $h$ values of $m$, independently of
the length of that gap. This bounded multiplicity is the step that
turns the interval construction into a count of consecutive gaps.

The weights are squares of signed sums of smooth divisor sums, each
built from divisors of shifted integers in the second interval.
The test functions make terms involving different shift sets
orthogonal in the limiting unmarked second moment.
At a prime belonging to one of those shift sets, the support cutoff
forces the corresponding divisor coordinate to equal one in every
nonzero summand. Removing that shift then couples terms whose
sets differ in size by one. The central choice is an alternating
family for which these neighboring terms nearly cancel.
The resulting square has positive average but arbitrarily small
weighted prime mass in the second interval. A detector estimate
for the first interval has a constant independent of the size of
the chosen family, which allows the suppression to be made strong
enough before the final counting argument.

The analytic ingredients are smooth divisor-sum correlations, the
Bombieri--Vinogradov distribution theorem, and an average of the
prime-tuple singular series. The distribution theorem originates
in the work of Bombieri and A.~I.~Vinogradov
\cite{Bombieri1965,Vinogradov1965}.
Goldston--Y{\i}ld{\i}r{\i}m's correlation method is a direct
antecedent: they square a linear combination of tuple approximations
of different cardinalities
\cite[p.~6, equation~(2.14) and the preceding paragraph]{GY2005}.
Maynard's smooth multidimensional sieve and product-profile
construction \cite{Maynard2015} provide further antecedents.
Here the orthogonality and coordinate-removal identities reduce the
choice of a weight to cancellation between adjacent dimensions.
All dimensions and test functions are fixed before $X$ tends to
infinity; the uniform error estimates needed to sum over the shifts
are proved in the paper.

Section~\ref{sec:counting} first deduces the theorem from a precise
weight proposition and proves the passage to ordinary index density.
Section~\ref{sec:weights} defines the square and states its moment
identities; Section~\ref{sec:cancellation} chooses the test functions
that make the adjacent dimensions cancel.
Section~\ref{sec:moments} proves the general mixed-moment estimates,
the singular-series average, and the moment identities used in the
construction.

\section{From two adjacent intervals to gap counts}\label{sec:counting}

For an integer $X\ge2$ and a fixed $\lambda>0$, write
\[
 L=\log X,\qquad h=\lfloor\lambda L\rfloor,\qquad
 J=\{1,\ldots,h\},\qquad I=\{h+1,\ldots,2h\}.
\]
We use the uniform average
$\mathbb E_X A=X^{-1}\sum_{X<m\le2X}A(m)$.
Set $\vartheta(n)=(\log n)\mathbf1_{\{n\ {\rm prime}\}}$ and
\[
 V_B(m)=\frac1L\sum_{b\in B}\vartheta(m+b)\qquad(B=I,J).
\]
The condition $V_B>0$ says exactly that the corresponding interval
contains a prime. Moreover, $V_B\ge1$ whenever $V_B>0$, since all
its prime arguments exceed $X$.

\begin{proposition}[Weights for adjacent intervals]\label{prop:weights}
For each fixed $\lambda>0$ there is a finite constant $K_\lambda>0$
with the following property. For every $\varepsilon>0$, there are
weights $W_X(m)\ge0$, defined for integers $X<m\le2X$, such that
\begin{align}
 \mathbb E_X W_X&\longrightarrow1,
 &\mathbb E_X(W_XV_J)&\longrightarrow\lambda,\label{eq:weight-first}\\
 \limsup_{X\to\infty}\mathbb E_X(W_XV_I)&\le\varepsilon,
 &\limsup_{X\to\infty}\mathbb E_X(W_XV_J^2)&\le K_\lambda,\label{eq:weight-primes}\\
 \mathbb E_X W_X^2&=O_{\lambda,\varepsilon}(1).\label{eq:weight-second}
\end{align}
The constant $K_\lambda$ is independent of $\varepsilon$.
\end{proposition}

The construction and proof occupy Sections~\ref{sec:weights}
through~\ref{sec:moments}. We now show why precisely these estimates
suffice. In particular, the bound in \eqref{eq:weight-second} need
not be uniform as $\varepsilon$ decreases.

The estimate for $\mathbb E_X(W_XV_J^2)$ prevents the detected prime
mass from being concentrated on too little weighted mass. After removing
the starts that contain a prime in $I$, the separate estimate for
$\mathbb E_XW_X^2$ prevents the remaining weighted mass from being
concentrated on too few ordinary starts. This is why only the first
constant must be chosen before the suppression tolerance.

\begin{lemma}\label{lem:adjacent-gaps}
Suppose $h\ge1$ and that for at least $\delta X$ integers $m\in(X,2X]$, the interval
$(m,m+h]$ contains a prime and $(m+h,m+2h]$ contains none.
Then at least $\delta X/h$ distinct primes $p\in(X,2X+h]$ have a
consecutive successor $p^+$ satisfying $p^+-p>h$.
\end{lemma}
\begin{proof}
For each such $m$, let $p(m)$ be the last prime in $(m,m+h]$.
There is no prime in $(p(m),m+2h]$, so
$p(m)^+>m+2h\ge p(m)+h$.
For a fixed selected prime $p$, the condition $m<p\le m+h$ implies
$p-h\le m\le p-1$. This interval contains exactly $h$ integers.
Thus each prime is selected at most $h$ times, proving the claim.
\end{proof}

\begin{figure}[ht]
\centering
\setlength{\unitlength}{1pt}
\begin{picture}(420,91)
 \put(10,36){\line(1,0){390}}
 \put(10,32){\line(0,1){8}}
 \put(205,32){\line(0,1){8}}
 \put(400,32){\line(0,1){8}}
 \put(10,16){\makebox(0,0){$m$}}
 \put(205,16){\makebox(0,0){$m+h$}}
 \put(400,16){\makebox(0,0){$m+2h$}}
 \put(95,49){\makebox(0,0){contains a prime}}
 \put(305,49){\makebox(0,0){no primes}}
 \put(148,36){\circle*{5}}
 \put(148,23){\makebox(0,0){$p(m)$}}
 \put(148,67){\vector(1,0){267}}
 \put(283,81){\makebox(0,0){no next prime before $m+2h$}}
\end{picture}
\caption{The last prime in the first interval begins a gap longer than
$h$. Its distance from $m$ is at most $h$, which bounds the counting
multiplicity even when the next prime is much farther away.}
\label{fig:blocks}
\end{figure}

\begin{proof}[Proof of Theorem~\ref{thm:main}]
Fix $C>0$ and choose $\lambda>\max\{C,1\}$.
Take $\varepsilon=\lambda^2/(2K_\lambda)$ in
Proposition~\ref{prop:weights}, and fix the resulting weight family.
There is a finite constant $M>0$ such that
$\mathbb E_XW_X^2\le M$ for every sufficiently large $X$.
Let $A_X=\{V_J>0\}$ and $B_X=\{V_I>0\}$.
Weighted Cauchy--Schwarz gives
\[
 \mathbb E_X(W_X\mathbf1_{A_X})
 \ge \frac{\bigl(\mathbb E_X(W_XV_J)\bigr)^2}
             {\mathbb E_X(W_XV_J^2)}
\]
for all sufficiently large $X$; its denominator is positive because
the numerator tends to $\lambda^2>0$.
By \eqref{eq:weight-first}--\eqref{eq:weight-primes},
\[
 \liminf_{X\to\infty}\mathbb E_X(W_X\mathbf1_{A_X})
 \ge\frac{\lambda^2}{K_\lambda}.
\]
Since $\mathbf1_{B_X}\le V_I$ and $W_X\ge0$, it follows that
\[
 \liminf_{X\to\infty}\mathbb E_X
      (W_X\mathbf1_{A_X\setminus B_X})
 \ge \frac{\lambda^2}{K_\lambda}-\varepsilon
 =:\Delta>0.
\]
Another use of Cauchy--Schwarz yields, for every sufficiently large $X$,
\[
 \frac{|A_X\setminus B_X|}{X}
 \ge\frac{\bigl(\mathbb E_X
             (W_X\mathbf1_{A_X\setminus B_X})\bigr)^2}
           {\mathbb E_XW_X^2}
 \ge\frac{\Delta^2}{4M}=:\delta>0.
\]
This is an ordinary proportion of integer starting points.
Lemma~\ref{lem:adjacent-gaps} therefore gives at least
\[
 \frac{\delta X}{h}\ge\frac{\delta X}{\lambda\log X}
\]
distinct prime starts $p\in(X,2X+h]$ with consecutive gap exceeding $h$.
Uniformly on this range,
$\log p=\log X+O(1)$ and $h=\lambda\log X+O(1)$.
The strict inequality $\lambda>C$ consequently gives
$h>C\log p$ for every sufficiently large $X$.

For an arbitrary sufficiently large integer $N$, choose
$X=\lfloor p_N/3\rfloor$. Then $2X+h<p_N$, so all these prime starts
have indices at most $N$. The prime number theorem gives
\[
 \frac{X}{\log X}\sim\frac{p_N}{3\log p_N}\sim\frac N3.
\]
In particular $X/\log X\ge N/4$ for every sufficiently large $N$.
The required count is at least $\delta N/(4\lambda)$.
Thus $c(C)=\delta/(4\lambda)>0$ is admissible.
All choices preceding $X$ depend only on the fixed threshold.
\end{proof}

The rest of the paper proves Proposition~\ref{prop:weights}.
The task is to make the weighted prime mass in $I$ arbitrarily small
without losing either the mass of the square or the uniform detector
bound in the adjacent block $J$.

\section{A square with a small marked moment}
\label{sec:weights}

We construct the weight in Proposition~\ref{prop:weights} as the square
of a signed sum of smooth divisor sums. We first define that square
and state the exact analytic estimates it satisfies. The next section
uses their formulas to choose the coefficients and functions; the
proofs of the analytic estimates follow in Section~\ref{sec:moments}.

Fix $\lambda>0$, and write
\[
 L=\log X,\qquad h=\lfloor\lambda L\rfloor,\qquad
 J=\{1,\ldots,h\},\qquad I=\{h+1,\ldots,2h\}.
\]
Here $X$ tends to infinity through integers, and $\mathbb E_X$ denotes
the average over the integers $X<m\le2X$.  Recall that
\[
 V_B(m)=\frac1L\sum_{b\in B}\vartheta(m+b),\qquad
 \vartheta(n)=(\log n)1_{\{n\ {\text{prime}}\}}.
\]
Let $\mu$ denote the M\"obius function. For a compactly supported
function $F:[0,\infty)^j\to\mathbb R$ and a tuple of distinct shifts
$\mathbf b=(b_1,\ldots,b_j)$, define
\[
 D_F(m;\mathbf b)
 =\sum_{d_i\mid m+b_i}\mu(d_1)\cdots\mu(d_j)
       F\left(\frac{\log d_1}{L},\ldots,
              \frac{\log d_j}{L}\right).
\]
The sum runs over positive divisors.
The zero-dimensional convention makes $D_F=F$ a scalar.  A symmetric
function $F$ allows us to write $D_F(m;S)$ for an unordered shift set $S$.

Put $\tau=1/8$, and let $k$ be a positive integer. For $1\le j\le k$, take a real symmetric function
\[
 f_j\in C_c^\infty
 \bigl(\{\mathbf t\in(0,\infty)^j:\textstyle\sum_i t_i<\tau\}\bigr).
\]
Allow also a real scalar $f_0$, and set $f_{k+1}=F_{k+1}=0$.  Define
\[
 F_j(\mathbf v)=\int_{s_i\ge v_i\ (1\le i\le j)}
                         f_j(\mathbf s)\,d\mathbf s,
 \qquad F_0=f_0,
 \qquad
 (Tf_{j+1})(\mathbf t)=\int_0^\infty f_{j+1}(\mathbf t,u)\,du.
\]
In dimension zero the last expression is a scalar.  All norms below
are Lebesgue $L^2$ norms on nonnegative orthants, with absolute value
in dimension zero. On the nonnegative orthant, $F_j$ vanishes when
$\sum_i v_i\ge\tau$. Cumulative integration also gives
\begin{equation}
 (-1)^j\partial_1\cdots\partial_jF_j=f_j.
 \label{eq:cumulative-derivative}
\end{equation}
Thus these complete mixed derivatives vanish on every coordinate
face, even though the cumulative functions themselves need not.

Define the signed sum
\[
 Z(m)=\sum_{j=0}^k\sum_{\substack{S\subset I\\|S|=j}}D_{F_j}(m;S),
 \qquad \alpha_j=\frac{\lambda^j}{j!}.
\]
The weight will be a positive normalization of $Z^2$. The following
analytic statement gives its mass, its interaction with the two prime
counts, and the second-moment bounds for prime detection and removal
of the weight. Its explicit formulas identify the remaining choice
of the functions $f_j$.

\begin{proposition}[Moment identities and a detector bound]
\label{prop:weight-moments}
Hold $\lambda$, $k$, and the family $(f_j)_{0\le j\le k}$ fixed.
Then, as $X\to\infty$,
\begin{align}
 \mathbb E_XZ^2&\longrightarrow
       w:=\sum_{j=0}^k\alpha_j\|f_j\|_2^2,
       &\mathbb E_XZ^4&=O_{\lambda,k,\mathbf f}(1),
       \label{eq:weight-unmarked}\\
 \mathbb E_X(Z^2V_J)&\longrightarrow\lambda w,
       &\mathbb E_X(Z^2V_I)&\longrightarrow
       v:=\lambda\sum_{j=0}^k\alpha_j
                         \|f_j+Tf_{j+1}\|_2^2.
       \label{eq:weight-marked}
\end{align}
Choose once and for all a real $G\in C_c^\infty(\mathbb R)$ with
$G(0)=1$ and $G(u)=0$ for $u\ge\tau$, and put
\[
 c_G=\int_0^\infty |G'(u)|^2\,du,
 \qquad C_*=\lambda+\lambda^2c_G^2.
\]
For every fixed family as above,
\begin{equation}
 \limsup_{X\to\infty}\mathbb E_X(Z^2V_J^2)\le C_*w.
 \label{eq:detector-bound}
\end{equation}
In particular, $C_*$ is independent of $k$ and of the family $(f_j)$.
\end{proposition}

The formulas distinguish two effects of a prime mark. A prime in $J$
lies outside every subset used to build $Z$, whereas a prime in $I$
alters the sum itself. To see the latter point exactly, fix $a\in I$.
When $m+a$ is prime, every nonzero term from a subset containing $a$
has divisor one in that coordinate: the other possible divisor satisfies
$\log(m+a)/L>1>\tau$, so its contribution vanishes by the support
condition on $F_j$.
Put $\widetilde F_j=F_j+F_{j+1}(\cdot,0)$. Deleting $a$ from each
subset that contains it therefore gives the exact identity
\begin{equation}\label{eq:prime-deletion}
 Z(m)=\sum_{j=0}^k
       \sum_{\substack{U\subset I\setminus\{a\}\\|U|=j}}
       D_{\widetilde F_j}(m;U)
 \qquad\text{when $m+a$ is prime}.
\end{equation}
Each subset occurs once, so deletion introduces no dimension factor.
The signed complete mixed derivative of $\widetilde F_j$ is
$f_j+Tf_{j+1}$. This explains why the marked form $v$ couples neighboring
dimensions.

The analytic proof evaluates the products in these squares by pairing
equal shifts. Any unmatched shift sets a coordinate of a complete
mixed derivative to zero, so the face vanishing in
\eqref{eq:cumulative-derivative} removes its main term. Only matching
subsets remain in the limiting quadratic forms, producing the squared
norms in $w$ and $v$. The same proof controls the errors and establishes
the two different second-moment bounds.

For a family with $w>0$, the normalization $W_X=Z^2/w$ therefore has
all the properties in Proposition~\ref{prop:weights} except the
prescribed smallness of the prime mass in $I$. That last requirement
is exactly $v/w\le\varepsilon$. The next section constructs families
with $w\to1$ and $v\to0$, while the detector constant $C_*$ remains
unchanged.

\section{Cancellation between adjacent dimensions}\label{sec:cancellation}

We now choose the functions so that the adjacent-dimension terms in
$v$ almost cancel.  Product functions restricted to a simplex, with
the truncation controlled by moments of their squared density, also
occur in Maynard's sieve construction \cite[Section~7]{Maynard2015}.
Here we combine such functions at a short range of dimensions with
alternating signs.

\begin{lemma}[A one-dimensional profile]
\label{lem:small-first-moment}
For every $\lambda>0$ and $\beta>0$, there exists a nonnegative
$g\in C_c^\infty((0,\infty))$ such that
\[
 \int_0^\infty g(u)^2\,du=1,\qquad
 \int_0^\infty g(u)\,du=\sqrt\lambda,\qquad
 \mu_g:=\int_0^\infty u g(u)^2\,du<\beta.
\]
\end{lemma}

\begin{proof}
For $R\ge2$, choose a smooth $\phi_R$ on $(0,\infty)$, equal to $1/u$
on $[1,R]$, with $0\le\phi_R(u)\le1/u$ and support in $[1/2,2R]$.
Put
\[
 A_R=\int\phi_R,\qquad B_R=\int u\phi_R(u)^2\,du,
 \qquad N_R^2=\int\phi_R^2.
\]
Then $A_R\ge\log R$, $B_R\le\log(4R)$, and
$1/2\le N_R^2\le2$.  Set $g_0=\phi_R/N_R$,
$c=\lambda/(\int g_0)^2$, and $g(u)=c^{-1/2}g_0(u/c)$.
Changes of variable give exactly
\[
 \int g^2=1,\qquad \int g=\sqrt\lambda,
 \qquad
 \mu_g=\lambda\frac{B_R}{A_R^2}
       \le\lambda\frac{\log(4R)}{(\log R)^2}.
\]
Choose $R$ sufficiently large and then hold it fixed.
\end{proof}

The profile just constructed may have a very large support endpoint
and variance. We need only their finiteness for this one fixed
profile: the dimension will be chosen afterwards. The auxiliary
limit in the next proposition concerns the constants $w$ and $v$
attached to finite families, before any asymptotic in $X$ is taken.

\begin{proposition}[Cancellation between high dimensions]
\label{prop:dimension-cancellation}
For every fixed $\lambda>0$, there are families of the type used in
Proposition~\ref{prop:weight-moments}, with increasing finite maximum
dimension $k$, for which $w\to1$ and $v\to0$.
\end{proposition}

\begin{proof}
Fix $0<\beta<\gamma<\tau$ and choose $g$ from
Lemma~\ref{lem:small-first-moment}.  Write $d=\beta-\mu_g>0$.
Let $B$ be an upper endpoint of the support of $g$, and let $s_g^2$
be the variance of the probability density $g^2$.  All these
quantities are fixed and finite; they need not be bounded uniformly
in $\lambda$.

For a large integer $k$ and $0\le j\le k$, define
\[
 Q_j(\mathbf t)=k^{j/2}\prod_{i=1}^j g(kt_i),\qquad
 S_j=\sum_{i=1}^j t_i,
\]
with $Q_0=1$ and $S_0=0$.  Before imposing the simplex cutoff, these
products satisfy
\[
 \|Q_j\|_2=1,\qquad TQ_{j+1}=\sqrt{\lambda/k}\,Q_j\quad(0\le j<k).
\]
Since $\alpha_j/\alpha_{j+1}=(j+1)/\lambda$, alternating products
$(-1)^jQ_j/\sqrt{\alpha_j}$ at neighboring levels therefore cancel
up to the factor $1-\sqrt{(j+1)/k}$.  We use $r$ neighboring levels
near $k$: the condition $r/k\to0$ makes these factors small, while
$r\to\infty$ makes the two endpoint contributions small after
normalizing the squared mass equally across the levels.

Take $r=\lfloor\sqrt k\rfloor$ and $a=k-r+1$.  Choose once and for all a smooth cutoff
$\chi$ with $0\le\chi\le1$, equal to one on $[0,\beta]$ and zero on
$[\gamma,\infty)$.  At the $r$ indices $a\le j\le k$, set
\begin{equation}
 f_j(\mathbf t)=\frac{(-1)^j}{\sqrt r\sqrt{\alpha_j}}
                     Q_j(\mathbf t)\chi(S_j),
 \label{eq:active-levels}
\end{equation}
and set all other $f_j$, including $f_0$, to zero.  Each selected finite
family has compact interior support: every coordinate is bounded
away from zero by a positive support endpoint of $g$ divided by $k$,
and the cutoff forces $S_j\le\gamma<\tau$.

The density $Q_j^2$ has total mass one.  Under it the variables $kt_i$
are independent with density $g^2$, so, for $j\le k$,
\[
 \mathbb E_jS_j=\frac{j\mu_g}{k}\le\mu_g,
 \qquad \operatorname{Var}_j(S_j)=\frac{js_g^2}{k^2}
                                      \le\frac{s_g^2}{k}.
\]
Chebyshev's inequality therefore gives
\begin{equation}
 1-\frac{s_g^2}{kd^2}
 \le w=\frac1r\sum_{j=a}^k\mathbb E_j[\chi(S_j)^2]\le1.
 \label{eq:w-concentration}
\end{equation}
This proves $w\to1$.

For $a\le j\le k-1$, substituting $u=kt_{j+1}$ and using
$\alpha_j/\alpha_{j+1}=(j+1)/\lambda$ gives
\begin{align}
 \sqrt{\alpha_j}(f_j+Tf_{j+1})(\mathbf t)
 &=\frac{(-1)^jQ_j(\mathbf t)}{\sqrt r}
       \left[\chi(S_j)-\sqrt{\frac{j+1}{\lambda k}}
                  \int_0^\infty g(u)\chi(S_j+u/k)\,du\right].
 \label{eq:interior-cancellation}
\end{align}
For $k\ge2B/d$, Chebyshev's inequality gives, uniformly in $j\le k$,
\[
 \mathbb P_j(S_j>\beta-B/k)\le\frac{4s_g^2}{kd^2}.
\]
On the complementary event every cutoff in
\eqref{eq:interior-cancellation} equals one.  The bracket is then
$1-\sqrt{(j+1)/k}$, whose square is at most $(r/k)^2$.  Everywhere,
both terms in the bracket lie in $[0,1]$, since $g\ge0$,
$\int g=\sqrt\lambda$, and $j+1\le k$; its absolute value is at most
one.  After summing the $r-1$ interior terms, their total contribution
to $v/\lambda$ is at most
\[
 (r/k)^2+\frac{4s_g^2}{kd^2}.
\]

Here the probability estimate is with respect to $Q_j^2$, while the
coordinate-deletion integral contains $g(u)\,du$. The passage between
them uses no comparison of these measures: on the complementary event,
$S_j+u/k\le\beta$ holds for every $u$ in the support of $g$.
Thus the whole integral has its cutoff equal to one pointwise.

Exactly two further indices can contribute.  At the top, $j=k$, the
term is $\alpha_k\|f_k\|_2^2\le1/r$.  At the bottom,
$j=a-1=k-r$, only $Tf_a$ occurs, and
\[
 \sqrt{\alpha_{a-1}}Tf_a
 =\frac{(-1)^aQ_{a-1}}{\sqrt r}
      \sqrt{\frac a{\lambda k}}
      \int_0^\infty g(u)\chi(S_{a-1}+u/k)\,du.
\]
Its squared norm is at most $a/(kr)\le1/r$.  Consequently
\begin{equation}
 0\le\frac v\lambda
 \le\frac2r+\left(\frac rk\right)^2+\frac{4s_g^2}{kd^2}
 \longrightarrow0.
 \label{eq:v-cancellation-bound}
\end{equation}
Together with \eqref{eq:w-concentration}, this proves the proposition.
\end{proof}

\begin{proof}[Proof of Proposition~\ref{prop:weights}]
Fix $\lambda>0$ and one function $G$ as in
Proposition~\ref{prop:weight-moments}. Set
\[
 K_\lambda=C_*=\lambda+\lambda^2c_G^2.
\]
This choice is independent of $\varepsilon$. Given $\varepsilon>0$,
choose $\beta$, $\gamma$, $g$, and $\chi$ as in the preceding construction.
By Proposition~\ref{prop:dimension-cancellation}, one can select a
\emph{single finite} $k$ and its family with $w>0$ and
$v/w\le\varepsilon$.  Freeze this family and define
\[
 W_X(m)=\frac{Z(m)^2}{w}.
\]
All conclusions follow from Proposition~\ref{prop:weight-moments} by
division by $w$, or by $w^2$ for the second moment.  Once the construction
is chosen in terms of $\lambda$ and $\varepsilon$, the fixed-family
second-moment bound may be written $O_{\lambda,\varepsilon}(1)$.

The parameter order is essential: the auxiliary limit $k\to\infty$
produces one family, and only then does $X\to\infty$.  No divisor-sum
asymptotic uniform in $k$ is required.  For an arbitrary fixed gap
threshold $C>0$, the application in Section~\ref{sec:counting} chooses a fixed
$\lambda>\max\{1,C\}$ before making these choices.
\end{proof}

\section{Analytic estimates for the square}\label{sec:moments}

It remains to prove the moment identities and detector bound in
Proposition~\ref{prop:weight-moments}. The products in $Z^2$ contain two
divisor-sum factors, those in $Z^4$ contain four, and the detector estimate
uses six. We first evaluate all these products in one formula that allows
repeated shifts and one prime mark. We then average the resulting singular
series over the shifts in $I$ and $J$, and apply the formulas to the square.
All dimensions, support budgets and functions remain fixed as $X$ tends
to infinity.

The leading coefficient separates into a test-function constant, determined
by which coordinates share a shift, and a singular series carrying the
numerical shifts. This separation is what permits the subset expansions
of $Z$ to be summed uniformly.

For a finite set $\mathcal H$ of distinct integers, put
\[
 \nu_p(\mathcal H)=|\mathcal H\bmod p|,\qquad
 \mathfrak S(\mathcal H)=\prod_p
 (1-\nu_p(\mathcal H)/p)(1-1/p)^{-|\mathcal H|}.
\]
We set $\mathfrak S(\varnothing)=1$. For a nonempty set the product converges:
beyond its diameter, the factors are $1+O_{|\mathcal H|}(p^{-2})$. It is nonnegative
and may be zero. The product is one when $|\mathcal H|=1$.

\begin{proposition}[Uniform mixed moments]\label{prop:mixed-moments}
Put $L=\log X$, where $X$ tends to infinity through integers, and let
$\mathbb E_X$ denote the average over $X<m\le2X$. Fix nonnegative integers
$j_1,\ldots,j_R$. For each $r$, let $F_r$ be a real function on
$[0,\infty)^{j_r}$ admitting a smooth compactly supported extension to
$\mathbb R^{j_r}$, and suppose that on the orthant it vanishes outside
\[
 v_1+\cdots+v_{j_r}\le\rho_r.
\]
Here the support budgets $\rho_r\ge0$ are fixed. Dimension zero means a scalar,
with support budget zero. For positive divisors
define
\[
 D_{F_r}(m;\mathbf b^{(r)})=
 \sum_{d_{r,\ell}\mid m+b_{r,\ell}\ (1\le\ell\le j_r)}
 \prod_{\ell=1}^{j_r}\mu(d_{r,\ell})
 F_r\left(\left(\frac{\log d_{r,\ell}}L\right)_{\ell=1}^{j_r}\right).
\]
For $j_r=0$, this means $D_{F_r}=F_r$.

Let $\Gamma=\{(r,\ell):1\le r\le R,\ 1\le\ell\le j_r\}$, and fix a
surjection $\kappa:\Gamma\to\{1,\ldots,t\}$. Write
$\mathcal M_i=\kappa^{-1}(\{i\})$. Take distinct integer shifts
$a_1,\ldots,a_t$ with $|a_i|\le DL^B$, where $D>0$ and $B\ge1$ are fixed,
and set $b_{r,\ell}=a_{\kappa(r,\ell)}$. Choose $\delta\in\{0,1\}$.
In the marked case take one more
shift $a_0$, distinct from these and satisfying the same bound. Put
\[
 (\chi_\delta(m),\mathcal H)=
 \begin{cases}
 (1,\{a_1,\ldots,a_t\}),&\delta=0,\\
 (\vartheta(m+a_0),\{a_0,a_1,\ldots,a_t\}),&\delta=1,
 \end{cases}
 \qquad \vartheta(n)=(\log n)1_{\{n\text{ prime}\}}.
\]
Assume $\sum_r\rho_r<1$ if $\delta=0$, and $\sum_r\rho_r<1/2$ if
$\delta=1$. Then there is a real constant $\mathcal C$, depending only on
the functions and $\kappa$, such that
\begin{equation}\label{eq:mixed-moment}
 \mathbb E_X\left[\chi_\delta(m)\prod_rD_{F_r}(m;\mathbf b^{(r)})\right]
 =L^{-t}\left\{\mathfrak S(\mathcal H)\mathcal C
 +O\!\left(L^{-1/2}(\log L)^A\right)\right\}.
\end{equation}
Here $A$ is a fixed finite number depending only on $\sum_rj_r$; the implied
constant may depend on all the fixed data, but is uniform in the actual
shifts. In particular the error is additive when $\mathfrak S(\mathcal H)=0$.

If each $|\mathcal M_i|$ is at most two, the constant is
\begin{equation}\label{eq:mixed-derivative-constant}
 \mathcal C=
 \int\prod_r\left((-1)^{j_r}\partial_1\cdots\partial_{j_r}F_r\right)
       \big((y_{\kappa(r,\ell)})_{\ell=1}^{j_r}\big)
       \prod_{i:\,|\mathcal M_i|=2}dy_i,
\end{equation}
where $y_i=0$ for $|\mathcal M_i|=1$, and each integrated variable ranges over
$[0,\infty)$. The zero-dimensional derivative is the original scalar.
When $\Gamma$ is empty, take $t=0$ and $\mathcal C=\prod_rF_r$.
\end{proposition}

The integer $t$ counts distinct numerical shifts, whereas
$\sum_r j_r$ counts divisor coordinates with repetitions. The power
$L^{-t}$ is determined by the former. Repetitions are recorded by the
fibers $\mathcal M_i$ and affect the constant $\mathcal C$. For two
divisor-sum factors whose individual shift tuples have distinct entries,
these fibers have size at most two. Products of four such factors
may have larger fibers; their estimates will use the general
finiteness assertion rather than the displayed derivative formula.
The marked formula retains the same power $L^{-t}$, consistent with
$\vartheta(m+a_0)$ having average tending to one. The additional factor
$L^{-1}$ in $V_B$ is supplied when the marked shifts are summed.

The scalar divisor-sum correlations of Goldston--Y{\i}ld{\i}r{\i}m
\cite{GoldstonYildirim2003Triple,GY2005} and Tao's smooth Fourier formulation,
described by Maynard \cite[Section~6, p.~404]{Maynard2015} motivate this calculation. We include its
proof because coupled smooth functions, arbitrary equality patterns and
an additive error uniform in logarithmic shifts are needed here.
The choice of smooth cutoffs was also informed by the analysis of
smoothing and high moments of divisor sums by
Granville--Koukoulopoulos--Maynard \cite{GKM2021}.  This is motivation
for the cutoff choice, not a source of the mixed-moment proposition
proved here.

\subsection{Residue classes and the prime mark}
Our objective is to replace a product of divisor sums, possibly weighted by one shifted prime, by an explicit finite density sum. We use the ordinary endpoint Bombieri--Vinogradov theorem in the form recorded by Bombieri, Friedlander and Iwaniec \cite[p.~206, equation~(1.4)]{BFI1986}; their $\psi$ notation is defined on p.~204. The additional divisor weights and endpoint shifts required here are handled below.

We retain the notation of Proposition~\ref{prop:mixed-moments} and put
$H_X=DL^B$. The two arithmetic lemmas below also allow $B\ge0$.
Negative shifts cause no difficulty, since all shifted arguments are
positive for sufficiently large $X$.
We write $\Lambda$ for the von Mangoldt function, $\varphi(q)$ for Euler's
totient, $\omega(q)$ for the number of distinct prime divisors of $q$,
$d(q)$ for the number of positive divisors, and $\zeta$ for the Riemann
zeta function.

\begin{lemma}[Weighted distribution with uniform shifted endpoints]\label{lem:arithmetic-1}

Assume the Bombieri--Vinogradov theorem: for every $A>0$ there is $b(A)>0$ such that
\begin{equation}
\sum_{q\le x^{1/2}(\log x)^{-b(A)}}
 \max_{(c,q)=1}\left|\psi(x;q,c)-\frac{x}{\varphi(q)}\right|
 \ll_A x(\log x)^{-A}.
\label{eq:arithmetic-BV}
\end{equation}
Here $\psi(x;q,c)=\sum_{n\le x,\ n\equiv c\pmod q}\Lambda(n)$, and $q=1$ is included. Fix $0\le\sigma<1/2$, $K\ge1$, and define
\[
E_X(q)=\max_{\substack{a\in\mathbb Z,\ |a|\le H_X\\(c,q)=1}}
\left|\sum_{\substack{X+a<n\le2X+a\\n\equiv c\pmod q}}\vartheta(n)
       -\frac{X}{\varphi(q)}\right|.
\]
For every $A>0$,
\begin{equation}
\sum_{q\le X^\sigma}\mu^2(q)K^{\omega(q)}E_X(q)
 \ll_{A,\sigma,K,D,B}X L^{-A}.
\label{eq:arithmetic-1}
\end{equation}
The conventions are $\omega(1)=0$, $K^0=1$, and $\varphi(1)=1$.

\end{lemma}
\begin{proof} At either endpoint $x=X,2X$, the range $q\le X^\sigma$ is contained in the range of \eqref{eq:arithmetic-BV} for every fixed logarithmic saving once $X$ is sufficiently large. Removing proper prime powers from a $\psi$-sum costs at most
\[
\sum_{p^j\le3X,\ j\ge2}\log p\ll X^{1/2}L^2
\]
in each residue class. Indeed there are $O(X^{1/2}L)$ pairs $p,j\ge2$, and each summand is $O(L)$. Translating the two endpoints through any $|a|\le H_X$ costs $O((H_X+1)L)$, uniformly in the residue class. Thus, writing $Q=\lfloor X^\sigma\rfloor$, for every $A_0>0$,
\begin{equation}
\sum_{q\le Q}E_X(q)
 \ll_{A_0}X L^{-A_0}
       +Q\{X^{1/2}L^2+(H_X+1)L\}
 \ll_{A_0,\sigma,D,B}X L^{-A_0}.
\label{eq:arithmetic-2}
\end{equation}
The last absorption uses the fixed strict inequality $\sigma<1/2$. The same bound controls the maximum over all shifts at once; no union bound over shifts is used.

An interval of length $X$ contains at most $X/q+1$ integers in one residue class. Also
\[
\frac q{\varphi(q)}=\prod_{p\mid q}\frac p{p-1}
 \le2^{\omega(q)}\le d(q).
\]
Since $q\le Q<X$ for large $X$, these two facts give
\begin{equation}
E_X(q)\ll \frac Xq\{L+d(q)\}.
\label{eq:arithmetic-3}
\end{equation}
For every fixed real $T\ge1$,
\begin{equation}
\sum_{q\le Q}\frac{\mu^2(q)T^{\omega(q)}}q
 \le\prod_{p\le Q}(1+T/p)\ll_T L^{c_T}
\label{eq:arithmetic-4}
\end{equation}
for a finite constant $c_T$. The elementary bound $\sum_{p\le y}1/p=O(\log\log(3y))$ suffices: for $y\ge2$, compare $1/p\le e p^{-1-1/\log y}$ with $\log\zeta(1+1/\log y)$, and bound $\zeta(1+u)\le1+1/u$. No prime-number theorem is needed for this estimate.

As squarefree $q$ satisfy $d(q)=2^{\omega(q)}$, \eqref{eq:arithmetic-3}--\eqref{eq:arithmetic-4} imply, for some finite $c_K$,
\[
\sum_{q\le Q}\mu^2(q)K^{2\omega(q)}E_X(q)
 \ll_K X L^{c_K}.
\]
Cauchy--Schwarz, applied to the nonnegative numbers $E_X(q)$, therefore gives
\[
\begin{split}
\sum_{q\le Q}\mu^2(q)K^{\omega(q)}E_X(q)
&\le
 \left(\sum_{q\le Q}E_X(q)\right)^{1/2}
 \left(\sum_{q\le Q}\mu^2(q)K^{2\omega(q)}E_X(q)\right)^{1/2}\\
&\ll X L^{(c_K-A_0)/2}.
\end{split}
\]
Choose $A_0\ge c_K+2A$. This proves \eqref{eq:arithmetic-1}. \end{proof}

\begin{lemma}[Exact density sum and arithmetic error]\label{lem:arithmetic-2}

Use the notation and support assumptions of
Proposition~\ref{prop:mixed-moments}; this lemma requires only bounded
functions $F_r$ and also allows $B\ge0$ in the shift bound
$H_X=DL^B$. Put $M=|\Gamma|$ and $\sigma=\sum_r\rho_r$.
As in the proposition, assume $\sigma<1$ when $\delta=0$ and
$\sigma<1/2$ when $\delta=1$.

For a squarefree divisor tuple $\mathbf d=(d_\gamma)_{\gamma\in\Gamma}$,
put $q=[d_\gamma:\gamma\in\Gamma]$, with $q=1$ when $M=0$.
Call the tuple compatible if, for every $p\mid q$, the shifts
$b_\gamma$ with $p\mid d_\gamma$ have one common residue modulo $p$.
When $\delta=1$, require this residue to differ from $a_0\bmod p$.
Set
\[
 \rho_\delta(\mathbf d)=
 \begin{cases}
  \displaystyle\prod_{p\mid q}(p-\delta)^{-1},
       &\mathbf d\text{ compatible},\\
  0,&\text{otherwise},
 \end{cases}
\]
and define the finite density sum
\begin{equation}
 \mathcal A_\delta=
 \sum_{\mathbf d\text{ squarefree}}
 \rho_\delta(\mathbf d)\prod_{\gamma\in\Gamma}\mu(d_\gamma)
 \prod_rF_r\left(
       \left(\frac{\log d_{r,\ell}}L\right)_{\ell=1}^{j_r}\right).
\label{eq:arithmetic-5}
\end{equation}
Empty products are one. For every $A>0$, uniformly in the allowed shifts,
\begin{equation}
 \mathbb E_X\left[\chi_\delta(m)
             \prod_rD_{F_r}(m;\mathbf b^{(r)})\right]
 =\mathcal A_\delta+O_A(L^{-A}).
\label{eq:arithmetic-6}
\end{equation}
The implied constant may depend on the fixed functions, dimensions,
support budgets, $D,B,A$. For $\delta=0$ and $M\ge1$, the error is
\begin{equation}
 O\bigl(X^{\sigma-1}(1+\log X)^{M-1}
          \prod_r\|F_r\|_\infty\bigr).
\label{eq:arithmetic-7}
\end{equation}

\end{lemma}
\begin{proof} Expand the product of divisor sums. M\"obius factors kill all nonsquarefree divisors. Whenever the product of the $F_r$'s is nonzero, the support hypothesis gives
\begin{equation}
\prod_{\ell=1}^{j_r}d_{r,\ell}\le X^{\rho_r}
\quad\text{for every }r,
\qquad
q\le\prod_\gamma d_\gamma\le X^\sigma.
\label{eq:arithmetic-8}
\end{equation}
This uses the budget of the sum of the coordinates in each factor, and does not impose a budget on the dimension.

For a squarefree tuple, the conditions $d_\gamma\mid m+b_\gamma$ prescribe at prime $p\mid q$ the residues $m\equiv-b_\gamma\pmod p$ for all selected coordinates. These congruences are inconsistent unless the selected shifts agree modulo every $p\mid q$. When they agree, the Chinese remainder theorem prescribes exactly one class $m\pmod q$. In the unmarked case this is precisely compatibility, and the count is $X/q+O(1)$, with an absolute constant in the $O(1)$.

When $M\ge1$, the number of positive integer tuples whose product is at most $Y\ge1$ is at most
\begin{equation}
Y\left(\sum_{d\le Y}\frac1d\right)^{M-1}
 \le Y(1+\log Y)^{M-1}.
\label{eq:arithmetic-9}
\end{equation}
Indeed, fix the first $M-1$ entries and bound the number of possible last entries by $Y$ divided by their product, then enlarge each preceding summation to $d\le Y$. Apply \eqref{eq:arithmetic-9} with $Y=X^\sigma$, multiply the $O(1)$ count error by $\prod_r\|F_r\|_\infty$, and divide by $X$. This proves \eqref{eq:arithmetic-7}, hence \eqref{eq:arithmetic-6} in the unmarked case. If $M=0$, the unmarked average and \eqref{eq:arithmetic-5} are both exactly $\prod_rF_r$.

For the marked case, put $n=m+a_0$. On a residue class satisfying these divisor congruences, the residue of $n\pmod p$ is $a_0-b_\gamma$ whenever $p\mid d_\gamma$. Consequently this class is reduced modulo $q$ exactly when the additional marked compatibility condition holds. If it is not reduced, a prime $n$ in that class would have to be a prime divisor of $q$, hence satisfy $n\le q$. But \eqref{eq:arithmetic-8} and $\sigma<1/2$ give
\[
n\ge X-H_X> X^\sigma\ge q
\]
for large $X$. The marked sum on a nonreduced class is therefore exactly zero, not an unestimated exceptional term. On a reduced class, the marked sum is
\[
\frac X{\varphi(q)}+O(E_X(q)).
\]
Since $q$ is squarefree, $\varphi(q)=\prod_{p\mid q}(p-1)$, proving the main density factor in \eqref{eq:arithmetic-5}.

For fixed squarefree $q$, each prime dividing $q$ must be assigned to a nonempty subset of the $M$ divisor coordinates. There are at most $(2^M-1)^{\omega(q)}\le (2^M)^{\omega(q)}$ choices if $M\ge1$; support and compatibility only reduce this count. The empty tuple, if $M=0$, has $q=1$ and is handled directly by the same bound with $K=1$. Thus the total normalized error is at most
\[
\frac{\prod_r\|F_r\|_\infty}{X}
 \sum_{q\le X^\sigma}\mu^2(q)(2^M)^{\omega(q)}E_X(q),
\]
which is $O_A(L^{-A})$ by Lemma~\ref{lem:arithmetic-1}. This proves the marked case, including all scalar factors and the wholly zero-dimensional case. \end{proof}

\subsection{Fourier evaluation of the density sum}

The residue-class calculation has controlled the prime-counting error.
It remains to evaluate its finite density sum, separating a constant
determined by the test functions from the singular series of the shifts.

\begin{proof}[Proof of Proposition~\ref{prop:mixed-moments}]
Let $M=|\Gamma|$. By Lemma~\ref{lem:arithmetic-2}, it remains to
evaluate $\mathcal A_\delta$ from \eqref{eq:arithmetic-5}: the arithmetic
error is $O_N(L^{-N})$ for every fixed $N$.
If $M=0$, that sum is exactly $\prod_rF_r$, and the assertion follows.
We henceforth assume $M>0$.

Choose a smooth compact extension $\widetilde F_r$ of each positive-dimensional
function and define its Fourier transform by
\[
 \Phi_r(\mathbf u)=(2\pi)^{-j_r}\int_{\mathbb R^{j_r}}
 e^{\sum_\ell v_\ell}\widetilde F_r(\mathbf v)
 e^{\mathrm i\mathbf u\cdot\mathbf v}\,d\mathbf v.
\]
The function $\Phi_r$ is Schwartz, and Fourier inversion gives exactly
\begin{equation}\label{eq:fourier-inversion}
 F_r(\mathbf v)=\int_{\mathbb R^{j_r}}\Phi_r(\mathbf u)
 e^{-\sum_\ell(1+\mathrm i u_\ell)v_\ell}\,d\mathbf u
 \qquad(\mathbf v\in[0,\infty)^{j_r}).
\end{equation}
For scalar factors use the scalar itself. Put $\Phi=\prod_r\Phi_r$,
$z_\gamma=(1+\mathrm i u_\gamma)/L$, and
$z_S=\sum_{\gamma\in S}z_\gamma$.

Insert \eqref{eq:fourier-inversion} into the density sum
\eqref{eq:arithmetic-5}, retaining the sum over all squarefree tuples.
The following bound justifies moving that sum inside the integrals:
\begin{equation}\label{eq:absolute-euler-sum}
 \sum_{\mathbf d}\rho_\delta(\mathbf d)\prod_\gamma d_\gamma^{-1/L}
 \le\prod_p(1+C_Mp^{-1-1/L})\ll_M L^{C_M}.
\end{equation}
Indeed $1/(p-\delta)\le2/p$, every nonempty selected subset contributes at
most $p^{-1/L}$, and there are at most $2^M-1$ subsets. Moreover
$\sum_pp^{-1-\varepsilon}\le\log\zeta(1+\varepsilon)
\le\log(1+1/\varepsilon)$. Multiplication by the finite integral
$\int|\Phi|$ proves absolute convergence of the sum-integral interchange.

Call a nonempty subset $S\subset\Gamma$ allowed at $p$ if its shifts satisfy
the same compatibility and mark-exclusion conditions. The exact Euler factor
is then
\[
 P_p(\mathbf z)=1+\frac1{p-\delta}
 \sum_{S\text{ allowed at }p}(-1)^{|S|}p^{-z_S},
 \qquad
 \mathcal A_\delta=\int\Phi(\mathbf u)\prod_pP_p(\mathbf z)\,d\mathbf u.
\]
By \eqref{eq:absolute-euler-sum} and Schwartz decay, the part of this integral
where $\max_\gamma|u_\gamma|>\sqrt L$ is $O_N(L^{-N})$ for every fixed $N$.
All subsequent comparisons are made on the complementary region.

The factors carrying the zeta poles are
\[
 B_p(\mathbf z)=\prod_{i=1}^t\prod_{\varnothing\ne S\subset\mathcal M_i}
 (1-p^{-1-z_S})^{-(-1)^{|S|}},\qquad
 H_p(\mathbf z)=P_p(\mathbf z)/B_p(\mathbf z).
\]
For $\Re z_\gamma\ge0$, all denominators in this expression have absolute
value at least $1/2$, and $|H_p(\mathbf z)|\le\exp(C_M/p)$. If distinct shifts
in $\mathcal H$ have distinct residues modulo $p$, the allowed subsets are
exactly the nonempty subsets of the individual fibers $\mathcal M_i$.
Matching the terms of order $p^{-1}$ gives
\begin{equation}\label{eq:generic-local-factor}
 H_p(\mathbf z)=1+O_M(p^{-2}),
\end{equation}
uniformly for $\Re z_\gamma\ge0$.

At $\mathbf z=0$, summing signs over the nonempty subsets of each occupied
allowed residue class gives
\[
 P_p(0)=\frac{p-\nu_p(\mathcal H)}{p-\delta},\qquad
 B_p(0)=(1-1/p)^t.
\]
It follows that $H_p(0)$ is the local factor of $\mathfrak S(\mathcal H)$.
This identity holds also when that factor is zero.

We now compare the product of the $H_p$'s additively with its value at zero.
Take $Y=\max(2,2DL^B)$. For $p\le Y$, differentiation along the segment from
zero to $\mathbf z$ gives
\[
 |H_p(\mathbf z)-H_p(0)|\ll_M L^{-1/2}\frac{\log p}{p}.
\]
For example, $|p^{-z_S}-1|\le|z_S|\log p$ on that segment, while the denominators
stay bounded away from zero. Each mixed partial product in the additive
telescoping identity is at most
\[
 \exp\left(C_M\sum_{p\le Y}1/p\right)\ll_M(\log Y)^{C'_M}.
\]
Using even the elementary bound
$\sum_{p\le Y}(\log p)/p\le\sum_{n\le Y}(\log n)/n=O((\log Y)^2)$, we obtain
an error $O(L^{-1/2}(\log L)^A)$ for the finite product. For $p>Y$ there are
no collisions of distinct shifts, so \eqref{eq:generic-local-factor} shows
that both remaining tail products are $1+O_M(1/Y)$. Consequently
\begin{equation}\label{eq:zero-safe-comparison}
 \prod_pH_p(\mathbf z)=\mathfrak S(\mathcal H)
  +O\!\left(L^{-1/2}(\log L)^A\right),\qquad
 \mathfrak S(\mathcal H)\ll(\log L)^A.
\end{equation}
This argument never divides by a possibly zero local singular-series factor.

It remains to evaluate the zeta factors. Absolute Euler products for
$\Re z_\gamma>0$ give
\[
 \prod_pB_p(\mathbf z)=\prod_{i=1}^t
 \prod_{\varnothing\ne S\subset\mathcal M_i}
 \zeta(1+z_S)^{(-1)^{|S|}}.
\]
Write $w_\gamma=1+\mathrm i u_\gamma$ and
\[
 K_i(\mathbf u)=\prod_{\varnothing\ne S\subset\mathcal M_i}
 \left(\sum_{\gamma\in S}w_\gamma\right)^{(-1)^{|S|+1}}.
\]
The pole expansion $\zeta(1+z)=z^{-1}(1+O(|z|))$ and the identity
$\sum_{\varnothing\ne S\subset\mathcal M_i}(-1)^{|S|}=-1$ imply
\begin{equation}\label{eq:zeta-kernel}
 \prod_pB_p(\mathbf z)
 =L^{-t}\prod_iK_i(\mathbf u)\{1+O_M(L^{-1/2})\}.
\end{equation}
Every denominator in $K_i$ has positive integer real part, so the product of
the $K_i$ has at most polynomial growth. Thus
\[
 \mathcal C=\int_{\mathbb R^M}\Phi(\mathbf u)\prod_iK_i(\mathbf u)\,d\mathbf u
\]
is absolutely convergent. Integrating \eqref{eq:zero-safe-comparison} and
\eqref{eq:zeta-kernel} against the Schwartz transform, then restoring the
negligible frequency tails, proves \eqref{eq:mixed-moment}. The constant is
real by conjugation under $\mathbf u\mapsto-\mathbf u$.

Finally, a singleton fiber contributes $w_\gamma$, and a double fiber
contributes
\[
 \frac{w_\gamma w_{\gamma'}}{w_\gamma+w_{\gamma'}}
 =w_\gamma w_{\gamma'}\int_0^\infty
 e^{-(w_\gamma+w_{\gamma'})y_i}\,dy_i.
\]
After taking absolute values, the new variables contribute
$e^{-2\sum_i y_i}$; the remaining factor is a Schwartz function times a fixed
polynomial. Fubini therefore applies. Recombining the Fourier integrals by
differentiating \eqref{eq:fourier-inversion} proves
\eqref{eq:mixed-derivative-constant}.
\end{proof}

\subsection{Averaging the singular series}\label{sec:singular-average}

We next average the numerical-shift factor in
Proposition~\ref{prop:mixed-moments}. The following form allows different
coordinates to range over either of the two blocks. It is a box version
of Gallagher's singular-series mean \cite[p.~5, equation~(3)]{Gallagher1976}.
We give a self-contained proof by truncating the Euler product;
no statement about the distribution of primes in individual short intervals
is assumed.

\begin{lemma}[Singular series in boxes]\label{lem:singular-average}
Fix an integer $s\ge0$ and a constant $K\ge1$. For each positive integer $h$,
let $I_1,\ldots,I_s$ be intervals of $h$ consecutive integers, all contained
in an interval of diameter at most $Kh$. Then, as $h\to\infty$,
\[
 \sum_{\substack{a_i\in I_i\ (1\le i\le s)\\a_1,\ldots,a_s\text{ distinct}}}
 \mathfrak S(\{a_1,\ldots,a_s\})=h^s(1+o(1)).
\]
The error is uniform in the intervals subject to this condition, with $s,K$
fixed. Here $A_s$ denotes a fixed finite constant depending only on $s$. In fact the normalized sum differs from one by
\[
 O_{s,K}(1/\log\log h)+O_s\big(h^{-1/2}(\log\log h)^{A_s}\big).
\]
\end{lemma}

\begin{proof}
The cases $s=0,1$ are immediate. Otherwise let $y=(1/4)\log h$ and
$Q=\prod_{p\le y}p$. The prime number theorem gives $Q\le h^{1/2}$ for large
$h$. For all tuples, including those with equal entries, define
\[
 \mathfrak S_y(\mathbf a)=\prod_{p\le y}
 (1-\nu_p(\{a_1,\ldots,a_s\})/p)(1-1/p)^{-s}.
\]
The exponent remains $-s$ even on tuples with collisions. For independent
uniform residues $a_1,\ldots,a_s$ modulo a prime $p$, averaging over an
additional uniform residue $x$ gives
\[
 \mathbb E(1-\nu_p/p)
 =\mathbb P(a_i\ne x\text{ for every }i)=(1-1/p)^s.
\]
The Chinese remainder theorem therefore shows that $\mathfrak S_y$ has
exactly mean one on complete residue tuples modulo $Q$.

A uniform integer in a length-$h$ interval has residue distribution modulo
$Q$ within total variation $O(Q/h)$ of uniform. The elementary bound $\sum_{p\le y}1/p=O(\log\log(3y))$ used above gives
\[
 0\le\mathfrak S_y(\mathbf a)
 \le\prod_{p\le y}(1-1/p)^{-s}\ll_s(\log y)^{A_s}.
\]
It follows that
\[
 h^{-s}\sum_{a_i\in I_i}\mathfrak S_y(\mathbf a)
 =1+O_s\big((Q/h)(\log y)^{A_s}\big).
\]
There are $O_s(h^{s-1})$ tuples with equal entries, so deleting them changes
this normalized mean by at most $O_s(h^{-1}(\log y)^{A_s})$.

For a distinct tuple put $\Delta=\prod_{i<j}|a_i-a_j|$, a positive integer.
For $y>2s$ every factor of the tail
\[
 T_y(\mathbf a)=\prod_{p>y}(1-\nu_p/p)(1-1/p)^{-s}
\]
is positive. A prime not dividing $\Delta$ has $\nu_p=s$ and a local factor
$1+O_s(p^{-2})$; the logarithm of a remaining factor is $O_s(1/p)$. Hence
\[
 |\log T_y|\ll_s \frac1y+\sum_{\substack{p>y\\p\mid\Delta}}\frac1p
 \le C_s\left(\frac1y+\frac{\log\Delta}{y\log y}\right).
\]
Since $\log\Delta\le\binom s2\log(Kh)$, this gives, uniformly in the distinct
tuples,
\[
 T_y(\mathbf a)=1+O_{s,K}(1/\log\log h).
\]
Now use the identity $\mathfrak S(\{a_1,\ldots,a_s\})=\mathfrak S_y(\mathbf a)T_y(\mathbf a)$
and the nonnegative truncated mean. This identity is valid also if a
small-prime factor vanishes; no quotient by $\mathfrak S_y$ is taken.
The stated error follows from $Q\le h^{1/2}$ and $\log y\asymp\log\log h$.
\end{proof}

When the shifts are summed, an equality pattern with $t$ distinct unmarked shifts has
$O(h^t)$ assignments in fixed blocks of length $h\asymp L$. Thus the uniform
error in Proposition~\ref{prop:mixed-moments}, summed over that pattern,
is $O(L^{-1/2}(\log L)^A)=o(1)$. If one additional distinct prime shift is
summed and its weight divided by $L$, there are $O(h^{t+1})$ assignments and
the same conclusion holds. Only finitely many patterns occur for a fixed
family of functions. The family is chosen before $X$ tends to infinity.

\subsection{Evaluation of the square}\label{sec:weight-moment-proof}

We now apply the correlation formula and singular-series mean to the
functions constructed in Section~\ref{sec:weights}. The signed complete
mixed derivative of each cumulative function is $f_j$, as in
\eqref{eq:cumulative-derivative}; its vanishing on coordinate faces will
leave only matching subsets in the quadratic moments.

The support restrictions in the moment formula apply to the sum of
the support budgets of the factors, independently of their dimensions.
For the present choice $\tau=1/8$, the required uses are as follows:
\begin{center}
\begin{tabular}{lll}
Quantity & Moment used & Total support budget\\
\hline
$\mathbb E_XZ^2$ & unmarked & $2\tau<1$\\
$\mathbb E_X(Z^2V_B)$, $B=I,J$ & one prime mark & $2\tau<1/2$\\
$\mathbb E_XZ^4$ & unmarked & $4\tau<1$\\
Detector off-diagonal terms & unmarked & $6\tau<1$
\end{tabular}
\end{center}
The last line uses divisor-sum majorants for the two possible primes
in $J$. It therefore requires no asymptotic formula with two prime marks.

\begin{proof}[Proof of Proposition~\ref{prop:weight-moments}]
First we check the smoothness and support needed by the moment
formula.  Extend $f_j$ by zero to $\mathbb R^j$.  Its cumulative integral
is smooth on $\mathbb R^j$, and multiplication by smooth coordinate
cutoffs that equal one near $[0,\infty)$ and vanish below $-1$ makes it
compactly supported without changing its orthant values or derivatives.
Indeed, the cumulative integral already vanishes when any coordinate
exceeds an upper support bound of $f_j$.  On the nonnegative orthant,
$F_j$ vanishes if $\sum_i v_i\ge\tau$, and its signed complete mixed derivative is given by
\eqref{eq:cumulative-derivative}.
This derivative is zero on every coordinate face.  For $j\ge1$,
$Tf_{j+1}$ has the same smoothness, interior-support, and face-vanishing
properties.

Expand $Z^2$ over pairs of subsets $S,T\subset I$.  Any shift in
$S\mathbin\triangle T$ occurs once and sets a coordinate of one of the
functions in \eqref{eq:cumulative-derivative} to zero.  Its main
coefficient is therefore zero.  If $S=T$ has size $j$, the coefficient
is $\|f_j\|_2^2$.  Averaging the singular series and using
\[
 \frac{\binom hj}{L^j}\longrightarrow\alpha_j
\]
proves the first limit in \eqref{eq:weight-unmarked}.  With a prime mark
in $J$, the same matching rule applies and
\[
 \frac{h\binom hj}{L^{j+1}}\longrightarrow\lambda\alpha_j.
\]
The two divisor factors have total support $2\tau<1/2$, so the marked
moment formula proves the first limit in \eqref{eq:weight-marked}.

For an internal mark $a\in I$, the exact deletion identity
\eqref{eq:prime-deletion} replaces $F_j$ by
$\widetilde F_j=F_j+F_{j+1}(\cdot,0)$ on the prime event.
Deletion is a bijection, so no extra dimension factor occurs.  The
support radius of $\widetilde F_j$ is still $\tau$, and its signed
complete mixed derivative is $f_j+Tf_{j+1}$.  It vanishes on coordinate
faces when $j\ge1$.  Apply the external marked formula to these
functions.  Only equal subsets $U$ survive, and
\[
 \frac{h\binom{h-1}{j}}{L^{j+1}}\longrightarrow\lambda\alpha_j.
\]
This proves the formula for $v$, including the scalar term
$\lambda|f_0+\int f_1|^2$.

The equality-pattern estimate preceding this subsection justifies
summing the errors in these limits. Applied to $Z^4$, the same argument
uses four divisor factors, total support
$4\tau<1$, and only finiteness of their pattern constants.  Bounded
averages of the singular series then give the second assertion in
\eqref{eq:weight-unmarked}.

It remains to prove a detector estimate whose constant is independent
of the chosen family.  Write $q_X=\log(2X+2h)/L=1+o(1)$.  At a
relevant prime $m+b$, we have $D_G(m;b)=1$.  Consequently, pointwise,
\[
 \frac{\vartheta(m+b)}L\le q_XD_G(m;b)^2,
 \qquad
 \left(\frac{\vartheta(m+b)}L\right)^2
       \le q_X\frac{\vartheta(m+b)}L.
\]
The first inequality holds at composites as well, since its left side
is zero.  Because $Z^2\ge0$, these inequalities give
\begin{align*}
 \mathbb E_X(Z^2V_J^2)
 &\le q_X\mathbb E_X(Z^2V_J)\\
 &\quad+q_X^2\sum_{\substack{b,c\in J\\b\ne c}}
       \mathbb E_X\bigl[Z^2 D_G(m;b)^2D_G(m;c)^2\bigr].
\end{align*}
An off-diagonal summand contains six divisor factors, with total
support $6\tau<1$.  The shifts $b,c$ are paired and lie outside $I$.
The two subsets from $Z^2$ must again be identical.  For their common
size $j$, the coefficient is $c_G^2\|f_j\|_2^2$.  The pair $(b,c)$ is
ordered, and
\[
 \frac{h(h-1)\binom hj}{L^{j+2}}
      \longrightarrow\lambda^2\alpha_j.
\]
Thus the off-diagonal sum tends to $\lambda^2c_G^2w$.
Together with \eqref{eq:weight-marked}, this proves
\eqref{eq:detector-bound}.  Although the convergence rates may depend
on the fixed family, the resulting $C_*$ does not.
\end{proof}

\bibliographystyle{plain}
{\raggedright\bibliography{references}}
\end{document}
